% Reconstruction of the missing A.66 arithmetic proof.
% Keep separate from the fixed-depth lifting audit until independently read.
\subsection{The variable-exponent arithmetic construction}

\begin{lemma}[The exact coefficient-prime image]
\label{lem:ghost-variable-exponent-image}
Let \(d\ge2\), let \(24\mid C\), and let
\(n=h!\), where \(h\ge2d^2\) and every prime divisor
of \(C\) is at most \(h\). Put
\[
 E(Z)=Cn\prod_{r=2}^d(Z^r-1),\qquad D=Cd!n^d.
\]
For a prime \(p\mid n\), write
\[
 v=v_p(n),\quad e=v_p(d),\quad
 a_p=\begin{cases}1,&p\ne2,\\2,&p=2,\end{cases}
 \qquad
 W_p=v_p(C)+(d+1)v+v_p(d!)+a_p+e.
\]
The continuous map
\[
 F_p(z)=\exp\bigl(E(z)\log z\bigr)
\]
is a bijection from
\(1+p^v+p^{v+a_p}\mathbb Z_p\) onto
\[
 F_p(1+p^v)\bigl(1+p^{W_p}\mathbb Z_p\bigr).
 \tag{VX.1}
\]
Every point of this domain satisfies
\[
 v_p(z-1)=v,\qquad v_p(E(z))=v_p(D),
 \qquad W_p\le2v_p(D).
 \tag{VX.2}
\]
The map preserves congruences modulo \(p^L\) between
points of the domain. For positive integers in the domain
it is the ordinary power \(z^{E(z)}\).
\end{lemma}
\begin{proof}
Set \(z=1+p^v u\). We have the exact identity
\[
 E(z)\log z=K u^d H(u),\qquad
 K=Cn\,d!\,p^{dv},
 \tag{VX.3}
\]
where
\[
 H(u)=
 \prod_{r=2}^d
 \frac{(1+p^v u)^r-1}{r p^v u}
 \cdot\frac{\log(1+p^v u)}{p^v u}.
\]
We prove coefficientwise that
\[
 H\in1+p^{e+1}\mathbb Z_p\langle u\rangle.
 \tag{VX.4}
\]
For a polynomial factor the coefficient of
\(u^{k-1}\), \(k\ge2\), has valuation at least
\(v(k-1)-v_p(k)\), using
\(\binom rk/r=\binom{r-1}{k-1}/k\).
When \(p\le d\), the factorial gives \(v\ge2d\);
thus this is at least \(e+1\), since
\(v_p(k)\le\lfloor\log_p d\rfloor\le d-1\).
When \(p>d\), all these denominators are units,
\(e=0\), and the bound is at least one.
For the logarithmic factor the same expression occurs
for every \(k\ge2\). Its minimum is at least \(v\)
for odd \(p\), and at least \(v-1\) for \(p=2\).
Indeed \(v_p(k)\le k-2\) for odd \(p\) and
\(v_2(k)\le k-1\); the only lower bound needed at
two uses \(v\ge2d\ge4\). These estimates also
tend to infinity with \(k\), proving convergence
in the indicated restricted power-series ring.
They prove (VX.4).

Write \(u=1+p^{a_p}w\). Expanding (VX.3), (VX.4) gives
\[
 u^dH(u)=H(1)+p^{a_p+e}
                   \bigl(\lambda w+pG(w)\bigr),
 \quad
 \lambda\in\mathbb Z_p^\times,
 \quad G\in\mathbb Z_p\langle w\rangle.
 \tag{VX.5}
\]
To check the nonlinear terms explicitly, use
\(v_p\binom dk\ge e-v_p(k)\). After division
by \(p^{a_p+e}\), their valuations are at least
\(a_p(k-1)-v_p(k)\ge1\) for \(k\ge2\).
The perturbation \(u^d(H(u)-1)\) contributes
coefficients divisible by \(p\) after this
division. The linear coefficient is
\(d/p^e\) modulo \(p\), hence is a unit.
For every \(t\in\mathbb Z_p\), the equation
\(\lambda w+pG(w)=t\) is equivalent to the
strict contraction
\(w=\lambda^{-1}(t-pG(w))\). An integral
restricted power series is one-Lipschitz on
\(\mathbb Z_p\), so this equation has exactly
one solution. Thus the logarithmic image is the
full coset modulo \(p^{W_p}\) claimed in (VX.1).
The exponential and logarithm are inverse isometries
on \(p\mathbb Z_p\) for odd \(p\), and on
\(4\mathbb Z_2\) at two; (VX.3) lies in these domains.
This proves (VX.1).

For the exact exponent valuation, the same factor
estimates without the logarithmic factor give
\(v_p(z^r-1)=v+v_p(r)\) for \(2\le r\le d\).
Summation gives \(v_p(E(z))=v_p(D)\).
Furthermore
\[
 2v_p(D)-W_p
   =v_p(C)+(d-1)v+v_p(d!)-a_p-e\ge0,
\]
because \(v_p(d!)\ge e\) and \(v\ge a_p\).
Finally, for two inputs congruent modulo \(p^L\),
the polynomial \(E\) and the logarithm both
preserve that congruence on this domain.
Their product and the exponential therefore do
as well. The logarithm identity for integral
powers proves the last assertion.
\end{proof}

\begin{theorem}[Reconstruction of SECTIONS A.66]
\label{thm:ghost-general-exponent-arithmetic}
Fix \(d\ge2\), a positive integer \(C\) divisible
by \(24\), and an integer \(J\ge2d^2\) at least
every prime divisor of \(C\). Put
\[
 n_j=(j+J)!,\qquad D_j=Cd!n_j^d,\qquad
 E_j(Q)=Cn_j\prod_{r=2}^d(Q^r-1).
 \tag{VX.6}
\]
There are three sequences of distinct primes
\(q_j,r_j,s_j>3\), disjoint from one another,
for which
\[
 A=3\prod_{j\ge0}\frac{q_j^{E_j(q_j)}}{s_j^{E_j(s_j)}},
 \qquad
 B=-2\prod_{j\ge0}\frac{r_j^{E_j(r_j)}}{s_j^{E_j(s_j)}}
 \tag{VX.7}
\]
are an ordinary completed unit-equation pair.
Their valuations at two and three are respectively
\((0,1)\), \((1,0)\), and their other nonzero
valuation pairs are exactly
\[
 \begin{array}{c|c}
 q_j&(E_j(q_j),0)\\
 r_j&(0,E_j(r_j))\\
 s_j&(-E_j(s_j),-E_j(s_j)).
 \end{array}
 \tag{VX.8}
\]
Both coordinates have infinite support. In the global
power completion, \(A/3\) and \(B/(-2)\) are
\(C\)-th powers. The construction is a recursion
by finite CRT conditions and specified Chebotarev
classes; all its simultaneous field conditions are
compatible.
\end{theorem}
\begin{proof}
We give the construction, including the field
intersection needed for every prime choice.
Keep pairwise coprime integral coefficients
\[
 a_j=3u_j^{24},\quad b_j=-2v_j^{24},\quad
 c_j=w_j^{24},
\]
where \(u_j,v_j,w_j\) are positive and have
disjoint prime supports away from two and three.
Set \(B(D)=(D+1)^4\) and
\(H_j(p)=2v_p(D_j)+j+3\). The two inductive
requirements are
\[
 \begin{gathered}
 a_jX^{D_j}+b_jY^{D_j}=c_jZ^{D_j}
 \quad\hbox{has a unit solution over every }\mathbb Z_p,\\
 v_p(a_j+b_j-c_j)\ge H_j(p)
 \quad\text{for }p\le B(D_j).
 \end{gathered}
 \tag{VX.9}
\]
The initial triple \((3,-2,1)\) satisfies both.

\emph{The three radical fields and their abelian intersections.}
For a current triple \((a,b,c)\), and \(24\mid n\),
write \(F_n=\mathbb Q(\mu_n)\). Each field
\[
 F_n((b/c)^{1/n}),\quad
 F_n((a/c)^{1/n}),\quad
 F_n((-a/b)^{1/n})
 \tag{VX.10}
\]
has maximal abelian subfield exactly \(F_n\).
Indeed the three radicands have squareclasses
\(-2,3,6\), whose square roots belong to \(F_n\).
None is an odd-prime power in \(F_n\): its valuation
at two or three is one, so an odd-prime radical has
that prime degree by the Eisenstein criterion over
the corresponding local field. Its real radical
field is not normal over \(\mathbb Q\), since its
other roots include nonreal roots. It cannot be a
subfield of an abelian extension of \(\mathbb Q\).
None is a fourth power in \(F_n\), because its
fourth root would give respectively
\(2^{1/4},3^{1/4},6^{1/4}\) there: the precise
factors are
\[
 (b/c)^{1/4}=(v_j/w_j)^6\zeta_8\,2^{1/4},
 \quad (a/c)^{1/4}=(u_j/w_j)^6\,3^{1/4},
\]
\[
 (-a/b)^{1/4}=(u_j/v_j)^6\,6^{1/4}/\sqrt2.
\]
Those real quartic fields are again nonnormal,
and irreducibility follows at two or three.
Consequently the cyclic Kummer kernel \(N\)
in each field (VX.10) has order exactly \(n/2\).
To see the degree assertion directly, a kernel
of order \(m\) would fix the \(m\)-th power
of the chosen radical, making the radicand an
\((n/m)\)-th power in \(F_n\); the preceding
exclusions force \(n/m\le2\), and the existing
square root gives equality.

Conjugation acts on \(N\) by the cyclotomic
character. All its multipliers are odd, and
complex conjugation has multiplier \(-1\), so
\([N,G]=N^2\). The fixed field of \(N^2\)
contains the splitting field of the corresponding
pure quartic, a nonabelian dihedral extension
of degree eight. Hence \(G/N^2\) is nonabelian.
Since \(N/N^2\) has order two and \(G/N\)
is abelian, \([G,G]=N\). This proves the
maximal-abelian assertion.
For every modulus \(M\) divisible by \(n\),
the intersection of any field (VX.10) with
\(\mathbb Q(\mu_M)\) is therefore exactly
\(F_n\). A CRT unit class that is one modulo
\(n\) defines an automorphism of the cyclotomic
field restricting to the identity on that
intersection. Pairing it with the identity on
(VX.10) is an actual element of the compositum's
fibre-product Galois group. Chebotarev
\cite[Theorem 8.31]{GhostMilneANT} thus supplies
infinitely many primes in that CRT class
splitting completely in the specified field.

\emph{Triangular local choices.}
At stage \(j\), abbreviate \(n=n_j\),
\(D=D_j\), \(E=E_j\), and \((a,b,c)=(a_j,b_j,c_j)\).
Choose \(P\ge B(D_{j+1})\) at least every
coefficient-support prime. Process all primes
\(p\le P\) in increasing order. Before the
process choose precisions \(L_p\) such that
\[
 L_p\ge H_{j+1}(p),\quad p^{L_p}>P,
 \quad L_p\ge v_p(n)+a_p\quad(p\mid n).
 \tag{VX.11}
\]
For each of the three future prime variables,
we impose
\[
 \begin{cases}
 z\in1+p^{v_p(n)}+p^{v_p(n)+a_p}\mathbb Z_p,
                                      &p\mid n,\\
 z\in\mathbb Z_p^\times,\quad
 z^r\not\equiv1\pmod p\ (2\le r\le d),
                                      &p\nmid n.
 \end{cases}
 \tag{VX.12}
\]
When \(p\nmid n\), we have \(p>j+J\ge2d^2\),
so there are nonzero residues outside the forbidden
set, whose size is at most \(\sum_{r=2}^d r<d^2\).
Before processing \(p\), the previous precisions
fix the three inputs modulo \(p-1\), and therefore
fix separately the exponent classes
\(e_i(p)=E(z_i)\bmod(p-1)\), \(i=a,b,c\).
This is because every prime-power divisor of
\(p-1\) is smaller than \(P\) and its prime has
already been processed.

At a prime \(\ell\mid p-1\), (VX.2) and (VX.12)
give, for every input consistent with the choices,
\[
 v_\ell(E(z_i))=v_\ell(D).
\]
For \(\ell\nmid n\) both valuations are zero:
the factors of \(C\) and \(d!\) lie in \(n\),
and all factors \(z_i^r-1\) are units by (VX.12).
Thus the three finite-field power maps have
the same exact image,
\[
 \gcd(e_i(p),p-1)=\gcd(D,p-1),\qquad
 (\mathbb F_p^\times)^{e_i(p)}
                  =(\mathbb F_p^\times)^D.
 \tag{VX.13}
\]
Every forbidden residue in (VX.12) maps to one
under each of these maps. Its order divides
some \(r\le d\), hence divides \(d!\), and
also divides \(p-1\); (VX.13) therefore makes
its order divide the relevant exponent.

At \(p\mid n\), the three local maps are
the common map of Lemma~\ref{lem:ghost-variable-exponent-image}.
Start all three inputs at \(1+p^{v_p(n)}\).
Their common output is a unit \(F_0\).
The sum-error is \((a+b-c)F_0\), of valuation
at least \(H_j(p)\ge W_p\), by (VX.9), since
\(p\le B(D)\). At least one of \(a,b,c\)
is a unit. Vary its output within the full
coset \(F_0+p^{W_p}\mathbb Z_p\) to cancel
the error, and use the lemma's inverse to
obtain the required input. Thus the exact
equation \(aF(z_a)+bF(z_b)=cF(z_c)\) is solved
without changing any shell in (VX.12).

At \(p\nmid n\), retain the already fixed
Teichmuller exponent \(e_i(p)\). On a permitted
residue class with Teichmuller representative
\(\omega\), define
\[
 F_{i,p}(z)=\omega^{e_i(p)}
       \exp\bigl(E(z)\log(z/\omega)\bigr).
 \tag{VX.14}
\]
This equals the ordinary power \(z^{E(z)}\)
at each positive integer satisfying the earlier
congruences. For \(z=\omega(1+pt)\),
\[
 \frac{\log(F_{i,p}(z)/\omega^{e_i(p)})}{p}
       =E(\omega)t+pG_i(t),
 \qquad G_i\in\mathbb Z_p\langle t\rangle.
 \tag{VX.15}
\]
The coefficient \(E(\omega)\) is a unit by (VX.12).
The remaining coefficients are integral and
divisible by \(p\): those of
\(\log(1+pt)/p\) in degree \(k\ge2\)
are \(p^{k-1}/k\), divisible by \(p\) because
\(p>2d^2\ge8\); the polynomial factor preserves
this divisibility. The contraction argument
from (VX.5) proves that (VX.14) is a bijection
onto \(\omega^{e_i(p)}(1+p\mathbb Z_p)\),
and that it preserves congruences.

For \(p\le B(D)\), choose a common value
\(\lambda\in(\mathbb F_p^\times)^D\) and
permitted preimages under all three power maps.
When the image has more than one element take
\(\lambda\ne1\); the forbidden residues all
map to one. When the image is \(\{1\}\),
all nonzero residues are preimages and the
bound on the forbidden set leaves a choice.
All three output balls in (VX.14) are then
\(\widetilde\lambda(1+p\mathbb Z_p)\).
The error at their common Teichmuller value
is divisible by \(p\), by (VX.9); varying an
output with unit coefficient cancels it.
For \(p>B(D)\), use a unit solution of the
first line of (VX.9). Multiply its three
coordinates by a common unit so that none
of their \(D\)-th powers is one modulo \(p\).
This is possible because
\((p-1)/\gcd(D,p-1)>3\), whereas there are
at most three forbidden scaling values.
The three resulting power values are in the
image (VX.13), and their preimages avoid every
forbidden residue. The bijections (VX.14)
give inputs attaining these exact local
values. This includes primes dividing a
coefficient, since (VX.9) supplies unit
coordinates there as well.

Approximate the three chosen local inputs
modulo \(p^{L_p}\). Congruence preservation
in (VX.1) and (VX.14) keeps the sum-error
divisible by \(p^{H_{j+1}(p)}\).
Their residues then determine the exponent
classes at later primes. Thus the triangular
procedure is compatible at every step.

\emph{Prime selection and the new invariant.}
Put \(M=\prod_{p\le P}p^{L_p}\). The three
resulting CRT classes are units and are one
modulo \(n\), by (VX.12). Apply the exact
fibre-product construction following (VX.10)
to choose distinct primes \(q,r,s>P\) in
these classes, splitting respectively in
the three fields (VX.10). Set
\[
 a'=a q^{E(q)},\qquad b'=b r^{E(r)},\qquad
 c'=c s^{E(s)}.
 \tag{VX.16}
\]
The support and twenty-fourth-power requirements
are preserved. At every \(p\le P\), the
new sum-error has the required precision.
Vary a variable with unit coefficient at
\((1,1,1)\); the inequality
\(H_{j+1}(p)>2v_p(D_{j+1})\) gives a unit
\(D_{j+1}\)-solution by one-variable Hensel.

At a new coefficient prime \(Q\in\{q,r,s\}\),
the imposed shells give
\(v_\ell(Q-1)=v_\ell(n)\) for every
\(\ell\mid n\). Every prime divisor of
\(D_{j+1}\) not dividing \(n\) is at most
\(j+J+1\le P\); its restriction (VX.12)
excludes \(Q\equiv1\pmod\ell\). Therefore
\[
              \gcd(D_{j+1},Q-1)=n.
 \tag{VX.17}
\]
The \(D_{j+1}\)-power and \(n\)-power images
in \(\mathbb F_Q^\times\) are exactly equal.
At \(q\), splitting in the first field
(VX.10) makes \(c/b\) an \(n\)-th power;
the new \(r\)- and \(s\)-factors are also
\(n\)-th powers, so \(c'/b'\) is a
\(D_{j+1}\)-th power modulo \(q\).
Set \(X=Z=1\) and solve
\(Y^{D_{j+1}}=(c'-a')/b'\) by simple-root
Hensel. At \(r\), set \(Y=Z=1\) and use
the split ratio \(c/a\). At \(s\), use
the ratio \(-b/a\), the inverse of the
third radicand. Each root is a unit and
each derivative is a unit, since the new
primes exceed \(P\ge B(D_{j+1})\).

Every other prime exceeds \(B(D_{j+1})\)
and all coefficients are units there.
For completeness, with \(g=\gcd(D_{j+1},p-1)\),
the Jacobi-sum count of nonzero solutions of
\(\alpha X^{D_{j+1}}+\beta Y^{D_{j+1}}=1\)
is at least
\[
 p-2-3(g-1)-(g-1)(g-2)\sqrt p>0.
 \tag{VX.18}
\]
Here the trivial character pair contributes
\(p-2\), the \(3(g-1)\) pairs with a
trivial character or trivial product have
absolute value one, and all remaining
Jacobi sums have absolute value \(\sqrt p\).
The inequality follows from
\(p>(D_{j+1}+1)^4\) and \(g\le D_{j+1}\).
The character identity follows by expanding
the two indicators of power images, so it
counts precisely the desired nonzero pairs: each power-image
indicator has a factor \(1/g\), and each pair of image
values has exactly \(g^2\) preimages, so these factors
cancel.
The usual Gauss-sum identity
\(J(\chi,\psi)=G(\chi)G(\psi)/G(\chi\psi)\)
and \(|G(\chi)|=\sqrt p\) prove the stated
bound when all three characters are nontrivial;
orthogonality gives the other terms.
The point lifts by its unit coordinate
derivatives. All places in (VX.9) have now
been checked, completing the induction.

\emph{The actual limits.}
Every new exponent is divisible by \(Cn_j\).
The three coefficient products and the
products after dividing all exponents by
\(C\) therefore converge in the global
power completion: modulo any fixed power
all sufficiently late factors are exactly
trivial. For each fixed finite prime \(p\),
its valuation in each coefficient becomes
constant, and all later factors are units.
Their exponents, divisible by growing
factorials, eventually annihilate every
finite quotient of \(\mathbb Z_p^\times\).
Consequently the coefficient sequences
converge in \(\mathbb Q_p^\times\) to
ordinary nonzero \(a_p,b_p,c_p\).
The precision in (VX.9) tends to infinity
and eventually applies to this \(p\),
so \(a_p+b_p=c_p\). Their ratios give
the ordinary point \(t_p=a_p/c_p\),
\(1-t_p=b_p/c_p\), with exactly the
localized completed classes (VX.7).
At the real place their signs are \((+,-)\),
represented by the ordinary pair \((3,-2)\).
Distinctness of all selected primes proves
(VX.8) and infinite support. Dividing the
convergent exponent products by \(C\)
proves the stated global \(C\)-th power
assertion. This completes the reconstructed
arithmetic proof of A.66.
\end{proof}
